\section{Alignment with the Common Domain Model}\label{ch:cdm}

This chapter discharges constitution~\S6 --- every right and obligation the
legal contract creates is a machine-readable unit --- and the~\S2 Clarity
commitment that behaviour is carried as declared data. It discharges them by a
single demonstration: the industry's Common Domain Model (FINOS CDM, version
6.0.0) maps onto the ledger's declared-data vocabulary. The direction of
authority is the constitution's and does not turn. The ledger's objects are
derived from first principles; the CDM is exhibited \emph{against} them, never
consulted as their source. Where the two models disagree, the disagreement is a
finding about the mapping --- recorded, and carried by a lineage convention ---
and never a reason to alter a ledger primitive.

\subsection{Unit maps to product; a bespoke instrument maps to the trade}

A ledger unit's declared terms map to a CDM product. The counterparty-agnostic
template --- the CDM \texttt{NonTransferableProduct}, whose \texttt{EconomicTerms}
carry the \texttt{Payout} legs, the effective and termination dates, and the
notional --- is the image of the unit's product-graph terms (Chapter~3). For a
listed, fungible instrument this template is the whole image: the same terms
serve every holder.

For a bespoke bilateral instrument the faithful image is not the template but
the CDM \texttt{Trade}. The reason is identity. In CDM the collateral terms live
on \texttt{Trade.collateral} (a \texttt{CollateralProvisions}), so two executions
sharing one \texttt{NonTransferableProduct} but differing in their collateral
terms are distinct \texttt{Trade}s. The ledger reaches the same conclusion from
the other direction and for the same reason: collateral eligibility and the
governing regime are declared terms of a distinct agreement unit
(Chapter~9), and a unit whose economics are inseparable from those terms is a
distinct unit. Unit identity therefore corresponds to \texttt{Trade} identity,
not to \texttt{NonTransferableProduct} identity --- collateral terms are part of
identity in both models. The one structural difference is placement: CDM
co-locates \texttt{CollateralProvisions} inside the \texttt{Trade}, whereas the
ledger factors the same terms into a separate agreement unit and relates the two
by projection. That is a finding about the mapping, carried by the
agreement-reference lineage of Chapter~9; it changes no primitive.

\subsection{Business event maps to transaction}

The CDM \texttt{BusinessEvent} --- a before-state, a set of
\texttt{PrimitiveInstruction}s, and the resulting \texttt{TradeState} --- maps to
the ledger's recorded \emph{transaction}, the atomic bundle that advances the
fold. The correspondence is structural on both sides. CDM composes a
\texttt{BusinessEvent} from primitive instructions (quantity change,
termination, transfer, reset, observation, and their siblings); the ledger
composes a transaction from its four kinds of change (moves, unit-state changes,
position-state changes, terms changes). The map from ledger transactions to CDM
business events has the shape the constitution already requires of the ledger's
own arrows, and it inherits their properties:

\begin{itemize}[nosep,leftmargin=1.4em]
  \item \textbf{total} --- every recorded transaction has a CDM image, because
        every transaction is one of the four declared kinds of change and each
        maps to a primitive instruction or a qualified composition of them;
  \item \textbf{deterministic} --- the image is a function of the record alone;
  \item \textbf{composing} --- a transaction's four change-kinds map
        component-wise to the event's primitive instructions, so composition of
        changes is composition of instructions;
  \item \textbf{conserving} --- the ledger's transaction conserves value per
        (unit, coordinate) by construction; the CDM event does not carry this
        law, and the ledger supplies it;
  \item \textbf{idempotent under the cause-derived key} --- a re-presented cause
        commits once, under the identifier of constitution~\S5; CDM has no native
        idempotence key, and the ledger's key is the lineage carried across the
        map.
\end{itemize}

The division of labour is exact and worth naming: \emph{the ledger keeps the
economic substance, the CDM carries the business meaning}. The conserved moves
and the state written into the three homes are the substance, guaranteed by the
Transaction Executor; the CDM qualification --- that this event \emph{is} a
partial termination, a reset, a variance fixing --- is the meaning, layered on
top of substance the ledger has already made unforgeable.

\paragraph{Variance swap.} Unit VS-1 (one-year OTC variance swap on IDX, 252
fixings, strike 400 variance points, 1{,}000 USD per point) maps to a product
whose \texttt{EconomicTerms} carry one \texttt{PerformancePayout}; its underlier
is the index IDX as an \texttt{Observable}, and the variance computation sits in
the payout's \texttt{returnTerms} as \texttt{VarianceReturnTerms} --- the field on
which the CDM equity-variance qualification depends. Each of the 252 fixings ---
the recorded IDX official closes --- is a CDM \texttt{Observation}, an observed
value on the record; the ledger's accrual firing, which writes the accrued
sum-of-squared-returns statistic into UnitStatus for the next firing to find,
maps to a \texttt{BusinessEvent} whose \texttt{Reset} lands in the trade's
\texttt{resetHistory}. Each fixing is thus a recorded observation and each
accrual a recorded business event, exactly as the two models agree they should
be. Final realised variance 441 points settles $1{,}000 \times (441-400) =
\USD{41{,}000}$ to the realised-variance receiver: one CDM transfer primitive,
one ledger owned-cash move, one conserved fact.

\paragraph{Total return swap.} Unit TRS-1 (notional \USD{10{,}000{,}000},
financing 4\% quarterly) maps to a two-legged product: a
\texttt{PerformancePayout} for the reference leg, whose underlier is the NAV of
wallet W-QIS as an \texttt{Observable}, and an \texttt{InterestRatePayout} for the
financing leg. The first reset is a CDM \texttt{BusinessEvent} carrying a
\texttt{Reset} (reset value: the stamped NAV observation \USD{10{,}300{,}000})
and a transfer; the ledger records it as one transaction with a single net move
of \USD{200{,}000}, payer to receiver. The alignment makes a reporting point the
netting would otherwise erase: the report derives from the \emph{retained
business event} --- total-return leg \USD{300{,}000} against financing leg
\USD{100{,}000}, both recoverable by recomputation from the recorded reset and
NAV observation --- never from the \USD{200{,}000} net move alone. Substance is
the net move on the record; meaning is the gross legs the business event retains.

\subsection{Eligibility as declared data restates the CDM collateral model}

The ratified collateral ruling makes eligibility and haircuts declared terms of
the collateral-agreement unit, checked at the single door and never encoded in a
type. This is the CDM collateral model restated in the ledger's primitives. CDM
carries eligibility as an \texttt{EligibleCollateral} schedule --- asset class,
currency, and a per-line valuation percentage after the shape of CSA
Paragraph~13 --- read as data, not as a subtype hierarchy of collateral kinds.
The ledger's line-level valuation default (Chapter~9) is the same choice:
collateral value is the sum of per-line market value times the declared
percentage, package-level valuation admissible only where an agreement declares
it. That the two models, derived independently, land on \emph{eligibility as
declared data valued line by line} is the alignment this section exhibits; it is
shown, not borrowed.

\subsection{Closed enumerations and the product graph}

Two further alignments serve later chapters. First, the CDM is built on closed
enumerations --- the eight-way \texttt{Payout} choice, the option-type and
day-count enumerations, the finite instruction set --- and closed choice types
are exactly a generator universe. The ledger's product-graph node sets are
closed at registration for the same reason (Chapter~3), and the CDM enumerations
seed the generator universe of Chapter~15, so that the property tests range over
the industry's own catalogue of product and event shapes, not an invented one.

Second, the product graph of Chapter~3 --- nodes the closed lifecycle states,
edges the guarded transitions that emit transaction templates --- is the natural
translation target for CDM lifecycle representation. A CDM \texttt{BusinessEvent}
is an edge traversal: the before-\texttt{TradeState} is the source node, the
after-\texttt{TradeState} the target, the primitive instructions the edge's
action. The CDM lifecycle and the ledger graph are the same object described
twice; the graph is where a CDM event round-trips without loss.

\subsection{Gaps}

Where CDM 6.0.0 has no native form for a situation the ledger must represent,
the mapping states the gap and names the lineage convention the ledger carries
instead. None is a defect in the ledger; each is a place the industry model has
not yet reached, recorded so that a later CDM version can be met without
surprise.

\begin{center}
\small
\begin{tabular}{@{}p{0.30\linewidth}p{0.28\linewidth}p{0.34\linewidth}@{}}
\toprule
\textbf{Situation} & \textbf{CDM 6.0.0 native form} & \textbf{Lineage convention carried} \\
\midrule
Securities-lending recall & None --- SL lifecycle events are not primitive
  instructions & Lent-plane transaction keyed to the loan agreement unit
  (Chapter~9): open $\to$ recall $\to$ return \\
Locate-style pre-trade confirmation & None --- CDM is a post-trade model & The
  locate is a declared obligation on the agreement unit, with deadline and
  discharge predicate; the pre-trade act itself is out of ledger scope \\
Collateral re-use / rehypothecation event & Collateral is modelled; re-use is
  not a first-class event & Posting of mass on the received ray with a
  mandatory source-agreement reference, from which per-agreement re-use projects \\
The coordinated cascade (one cause, many units) & None --- events are
  per-trade & The cause-derived key of constitution~\S5 ties the cascade's
  transactions together and makes them jointly idempotent \\
\bottomrule
\end{tabular}
\end{center}

No gap adds a mechanism: the recall rides the lent plane, the locate rides the
obligation-liveness discipline, re-use rides the source-agreement reference, and
the cascade rides the cause-derived identifier. The mapping is total on the
ledger's side precisely because the residue is absorbed by primitives the design
already carries.
